\bookchapter{Chapitre 44 : De 0 à 1 en pratique — le parcours complet d'un projet}

\chapterintro{La théorie est terminée, place à la pratique. Ce chapitre n'est pas un simple tutoriel à recopier : c'est une méthode de terrain pour apprendre « comment penser, comment demander, comment vérifier ». À travers trois projets de difficulté croissante, nous verrons tout le cheminement, du besoin au produit fini. Quel que soit l'outil que tu utilises, ce processus s'applique.}

\booksection{Projet 1 : un mini-site d'outils personnels (niveau débutant)}

\booksubsection{Réfléchir au besoin}

Objectif : créer une page web d'outils personnels qui présente l'accès à 3 ou 4 petits outils pratiques. C'est un projet purement front-end, sans back-end ni base de données. L'important n'est pas la complexité des outils, mais de vivre le parcours complet « du besoin au produit fini ».

Selon les cinq éléments du besoin : Who — un usage personnel. What — une page web présentant l'accès aux outils. Why — regrouper les outils du quotidien sur une seule page, faciles d'accès. How — un fichier unique HTML+CSS+JS, sans framework. Constraint — à réaliser en 30 minutes, sans déploiement.

\booksubsection{Donner les instructions}

En suivant le principe « Think Before Coding », ton instruction doit contenir assez d'informations pour que l'IA comprenne l'objectif, sans pour autant préciser chaque détail d'implémentation :

\begin{lstlisting}
Crée une page web personnelle d'outils avec ces exigences :
1. HTML en un seul fichier (CSS et JS en ligne), sans framework
2. Thème sombre, style moderne et minimaliste
3. Affiche 4 cartes d'accès aux outils :
   - Résumeur de texte
   - Compteur de mots
   - Sélecteur de couleur
   - Générateur de QR codes
4. Chaque carte avec icône, titre et description courte
5. Les cartes ont une animation au survol
6. Mise en page responsive, utilisable sur mobile
7. Une barre de titre en haut
Donne directement le fichier HTML complet.
\end{lstlisting}

\booksubsection{Réception et itération}

Quand l'IA fournit le code, fais la réception avec les trois niveaux de vérification itérative. Premier niveau, fonctionnel : le fichier s'ouvre-t-il dans le navigateur ? Les 4 cartes s'affichent-elles toutes ? Les animations au survol fonctionnent-elles ? Deuxième niveau, limites : la mise en page tient-elle sur un écran de téléphone ? Les cartes passent-elles à la ligne quand la fenêtre est très étroite ? Troisième niveau, architecture : le code est-il lisible ? Le CSS est-il regroupé dans une seule balise style ? Y a-t-il des redondances évidentes ?

Si tu veux aussi implémenter la fonction d'un outil (et pas seulement son accès), procède par étapes — ajoute la fonction d'un outil, vérifie, puis ajoute le suivant. C'est l'application concrète de la stratégie des « petits pas rapides ».

\booksection{Projet 2 : une application de comptabilité full-stack (niveau intermédiaire)}

\booksubsection{Réflexion sur le besoin et PRD}

Objectif : une application de comptabilité full-stack avec front-end, back-end et base de données. Ce projet est bien plus complexe que le premier et mobilise la méthode d'ingénierie des besoins du chapitre 40.

Commence par rédiger le PRD. Sur le modèle de la section 3 du chapitre 40, décris la vue d'ensemble du projet, la liste des fonctionnalités, la solution technique, le modèle de données et le plan de développement. Pour la technique, applique le principe « Simplicity First » : React côté front mais sans gestion d'état complexe type Redux, Node.js côté back mais sans microservices, SQLite plutôt que PostgreSQL pour la base (en phase de développement).

Le plan de développement se découpe en 5 étapes : étape 1, le squelette de l'API back-end (serveur Express + SQLite + routes CRUD) ; étape 2, les pages de base du front-end (liste des enregistrements + formulaire d'ajout) ; étape 3, la liaison front-back (appels d'API + affichage des données) ; étape 4, les catégories et le filtrage ; étape 5, le récapitulatif mensuel et de simples graphiques. À chaque étape, tu obtiens quelque chose d'utilisable.

\booksubsection{Créer le fichier mémoire du projet}

Avant de coder, crée d'abord le fichier mémoire du projet. C'est la pratique centrale de l'ingénierie du contexte : l'IA connaît le contexte du projet dès le début de chaque conversation.

\begin{lstlisting}
# Projet : app full-stack de suivi de dépenses
## Aperçu
App de suivi de dépenses personnelles avec enregistrement des revenus et dépenses, catégories et résumés mensuels.
## Stack technique
- Frontend : React + TypeScript + Vite
- Backend : Node.js + Express + TypeScript
- Base de données : SQLite (via better-sqlite3)
- Types partagés entre frontend et backend
## Conventions de l'API
- Chemin de base : /api
- Format de réponse : { success: boolean, data: any, error: string }
- Style RESTful
## Modèle de données
Record { id, amount, type, category, note, date }
## Conventions de code
- Mode strict TypeScript
- Indentation de 2 espaces
- Fonctions avec commentaire JSDoc
- Commentaires et messages de commit en français
\end{lstlisting}

\booksubsection{Implémenter pas à pas}

Conformément au plan de développement, réalise chaque étape dans une conversation dédiée. Voici un exemple d'instruction pour la première étape :

\begin{lstlisting}
J'ai besoin de créer le squelette de l'API backend.
Exigences :
1. Serveur Express, port 3001
2. Base de données SQLite, table records
3. Champs de la table records :
   id (TEXT PRIMARY KEY)
   amount (REAL)
   type (TEXT: 'income'|'expense')
   category (TEXT)
   note (TEXT)
   date (TEXT: YYYY-MM-DD)
4. Endpoints CRUD :
   GET /api/records - récupérer toutes les entrées
   POST /api/records - ajouter une entrée
   PUT /api/records/:id - mettre à jour une entrée
   DELETE /api/records/:id - supprimer une entrée
5. Middleware de gestion des erreurs
Donne la structure complète du code backend.
\end{lstlisting}

Note ce qui fait la force de cette instruction : elle décrit uniquement l'étape en cours (le squelette back-end), sans y mêler de besoins front-end ; elle liste clairement les champs et les points d'entrée (assez de contraintes pour l'IA), tout en laissant le détail de l'implémentation à l'IA ; elle suit le principe « Goal-Driven Execution » — dire « quoi faire » plutôt que « comment faire ».

Dès que l'IA fournit le code, vérifie aussitôt : lance le serveur et teste chaque point d'entrée avec curl ou Postman. Une fois tout confirmé, ouvre une nouvelle conversation pour l'étape 2 (les pages de base du front-end). Chaque étape suit la boucle « instruction - vérification - nouvelle conversation ».

\booksection{Projet 3 : une application augmentée par l'IA (niveau avancé)}

\booksubsection{Concevoir le projet}

Objectif : une application connectée à l'API d'un grand modèle — par exemple un « assistant d'apprentissage IA » : l'utilisateur saisit un contenu à étudier, et l'IA produit un résumé, relie les notions entre elles et recommande des ressources. La particularité de ce projet : ton application appelle elle-même l'IA, et tu dois gérer l'appel d'API, l'ingénierie des prompts, la gestion des erreurs et autres nouveaux défis.

La leçon la plus importante de ce projet : quand tu « utilises l'IA pour écrire le code qui appelle l'IA », l'ingénierie des prompts demande une attention toute particulière. Les prompts que ton application envoie à l'API du grand modèle doivent eux aussi suivre les principes vus dans ce guide — clarté, concision, orientation objectif.

\booksubsection{Les points techniques clés}

Ce projet soulève plusieurs points techniques que tu n'as pas rencontrés dans les deux premiers :

\begin{itemize}
  \item Gestion des clés API : une clé ne doit jamais être codée en dur, gère-la via des variables d'environnement ;
  \item Traitement des réponses en flux : les réponses des grands modèles arrivent généralement en streaming, à afficher mot à mot ;
  \item Maîtrise de la consommation de tokens : les grands modèles facturent au token, il faut contrôler la consommation au niveau de l'application ;
  \item Gestion des modèles de prompts : chaque fonction a son propre modèle de prompt, à organiser de façon structurée.
\end{itemize}

Tu peux confier tous ces points à l'IA, mais tu dois en comprendre le « pourquoi » : pourquoi une clé ne doit pas être codée en dur ? pourquoi des réponses en flux ? pourquoi contrôler la consommation de tokens ? C'est en comprenant le « pourquoi » que tu pourras juger si la solution de l'IA tient la route.

\booksection{La liste de contrôle universelle des projets}

Quelle que soit la taille du projet, voici une liste de contrôle universelle pour programmer avec l'IA. À chaque nouveau projet, parcoure-la pour ne rater aucune étape clé.

\begin{center}
\begin{tabularx}{\linewidth}{|X|X|X|}
\hline
\textbf{Phase} & \textbf{Quoi faire} & \textbf{Principe associé} \tabularnewline
\hline
1. Réfléchir au besoin & Passer en revue les cinq éléments du besoin, clarifier ce que l'on veut faire & Think Before Coding \tabularnewline
2. Découper le besoin & Découper le besoin en petites étapes vérifiables une à une & Simplicity First \tabularnewline
3. Rédiger le PRD & Le formaliser dans un document structuré, pour s'y référer & Goal-Driven \tabularnewline
4. Créer le fichier mémoire & Créer CLAUDE.md ou un fichier équivalent & Harness Engineering \tabularnewline
5. Implémenter pas à pas & Une conversation par étape : instruction, vérification, étape suivante & Surgical Changes \tabularnewline
6. Tester et vérifier & Tester le bon fonctionnement dès chaque étape terminée & Iterate \& Verify \tabularnewline
7. Relire le code & Contrôler l'architecture, la sécurité, la maintenabilité & Iterate \& Verify \tabularnewline
8. Itérer et améliorer & Affiner par itérations à partir des tests et de la relecture & Iterate \& Verify \tabularnewline
\hline
\end{tabularx}
\end{center}

\begin{center}\small\color{gray}La liste de contrôle universelle de la programmation IA\end{center}

\begin{tcbpractice}{Passe à l'action : termine le « projet 1 » cette semaine}
Prends le « mini-site d'outils » de la section 1 du chapitre 44 comme ton premier devoir de programmation IA : 5 minutes pour les cinq éléments du besoin, envoie l'instruction à l'IA (n'importe quel outil fera l'affaire), puis fais la réception avec les trois niveaux de vérification. Tu as 30 minutes — et voilà ta première réalisation en programmation IA. Une fois terminé, garde-la précieusement : c'est le trophée de ton passage de « lecteur » à « praticien ».

\end{tcbpractice}

\begin{tcbsummary}{L'essentiel du chapitre}
\begin{itemize}
  \item Projet 1 (débutant) : un site d'outils pur front-end — tout le parcours en 30 minutes
  \item Projet 2 (intermédiaire) : une application de comptabilité full-stack — PRD + fichier mémoire + 5 étapes successives, chacune « utilisable »
  \item Projet 3 (avancé) : un assistant d'apprentissage IA — écrire avec l'IA le code qui appelle l'IA, en soignant les prompts
  \item Les huit étapes du processus universel : réflexion sur le besoin découpage PRD fichier mémoire implémentation pas à pas tests relecture itération
  \item Chaque étape suit la boucle « instruction - vérification - nouvelle conversation »
\end{itemize}

\end{tcbsummary}
